\section{Approximating the Longest Tour in a Directed Graph}
To construct the overlap approximation subroutine used in
$2\frac{25}{42}$-Approximation algorithm, we analyze the
$\frac{38}{63}$-approximation algorithm for the directed Max TSP. For the
simplicity in the theoretical analysis, we assume throughout this section that
the number of vertices in the input graph is even. If the number of vertices is
odd, we can add a vertex that is connected to all other vertices with $0$-weight
edges.

The algorithm heavily relies on two standard polynomial-time graph subroutines:
\begin{enumerate}
    \item The Hungarian algorithm \cite{kuhn1955hungarian, edmondskarp1972, tomizawa1971techniques} for solving the
    maximum-weight cycle cover problem via its reduction to the assignment
    problem, running in $O(n^3)$ time.
    \item Edmonds' blossom algorithm \cite{edmonds1965paths, edmonds1965matching} for finding a
    maximum-weight (not necessarily perfect) matching in an undirected graph,
    also running in $O(n^3)$ time.
\end{enumerate}

\subsection{The Generic Directed Max TSP Algorithm}
The approximation algorithm consists of three distinct versions, all of which
share the following framework:

\begin{enumerate}
    \item \textbf{Cycle Cover:} Compute a maximum-weight cycle cover
    $\mathcal{C}$ for the graph $G$.
    \item \textbf{Path Extraction:} Use the cycle cover $\mathcal{C}$ to obtain
    a set $P$ of vertex-disjoint paths.
    \item \textbf{Tour Construction:} Construct a final tour by patching
    together the paths from $P$.
\end{enumerate}

The three versions of the algorithm differ strictly in their implementation of
Step 2. The first version follows the simple approach of deleting the smallest
edge in each cycle. The second and third approaches are significantly more
complicated, utilizing non-cycle edges in an attempt to extract a set of paths
with a greater total weight.
\subsection{Preliminaries}
Before comparing these three approaches, we must establish the following
definitions regarding the cycle cover $\mathcal{C}$:
\begin{itemize}
    \item A \textbf{2-cycle} is any cycle passing through exactly 2 vertices.
    \item A \textbf{$3^+$-cycle} is any cycle passing through 3 or more
    vertices.
    \item $W$ denotes the total weight of all edges in $\mathcal{C}$.
    \item $bW$ denotes the total weight of the edges that are the heavier (i.e.,
    larger-weight) edge in some 2-cycle of $\mathcal{C}$.
    \item $cW$ denotes the total weight of the edges that are the lighter (i.e.,
    smaller-weight) edge in some 2-cycle of $\mathcal{C}$.
\end{itemize}

By expressing the approximation bounds of the three versions as linear functions
of the term $(b - 2c)$, we can prove that the best of the three always
guarantees at least a $\frac{38}{63}$-approximation.

\subsection{Crucial Observation}
A crucial observation simplifies the mathematical analysis of this framework. To
bound the weight of the final produced tour, it is entirely sufficient to obtain
a lower bound on $w(P)$, which is the total weight of the paths extracted in
Step 2. Because the graph strictly contains non-negative edge weights, no matter
how the vertex-disjoint paths of $P$ are patched together in Step 3, the weight
of the resulting tour will always be at least $w(P)$. As the result, the
explicit operations of Step 3 can be safely ignored in the remainder of the
analysis.
